\documentclass[a4paper]{article}

%% Language and font encodings
\usepackage[T1]{fontenc}
\usepackage[utf8x]{inputenc}
\usepackage[english]{babel}

\usepackage[colorlinks=true, allcolors=blue]{hyperref}

\urlstyle{tt}
\newcommand{\email}[1]{\href{mailto:#1}{\tt{\nolinkurl{#1}}}}
\newcommand{\orcid}[1]{ORCID: \href{https://orcid.org/#1}{\tt{\nolinkurl{#1}}}}

\usepackage[sfdefault,lf]{carlito}
%% The 'lf' option for lining figures
%% The 'sfdefault' option to make the base font sans serif
\usepackage[parfill]{parskip}
\renewcommand*\oldstylenums[1]{\carlitoOsF #1}
\usepackage{fancyhdr}
\usepackage{natbib}
\usepackage{authblk}
\setlength{\headheight}{41pt}

%% Sets page size and margins
\usepackage[a4paper,top=3cm,bottom=2cm,left=3cm,right=3cm,marginparwidth=1.75cm]{geometry}

%% Useful packages
\usepackage{amsmath}
\usepackage{graphicx}
\usepackage{booktabs}
\usepackage{color}

\usepackage[colorinlistoftodos]{todonotes}

\newcommand{\cz}[1]{\hl{\textbf{cz:} #1}}




%\renewcommand{\headrulewidth}{0pt}
% \fancyhead[L]{Posted: \today}
% \fancyhead[R]{
% \includegraphics[width=4cm]{./figures/engrXiv_banner.png}
% }
\pagestyle{plain}
\title{Data-efficient and interpretable inverse materials design using a disentangled variational autoencoder}
\author[1, $\dag$, *]{Zulqarnain Khan}
\author[1, 2, $\dag$, *]{Cheng Zeng}
\author[2]{Nathan L. Post}
\affil[1]{Institute for Experiential AI, Northeastern University, Boston, MA 02115, United States}
\affil[2]{The Roux Institute, Northeastern University, Portland, ME 04101, United States}
\affil[$\dag$]{These authors contribute equally: Zulqarnain Khan, Cheng Zeng.}
\affil[*]{Corresponding author: \email{z.khan@northeastern.edu}, \email{c.zeng@northeastern.edu}}
\date{}

\begin{document}
\maketitle
\thispagestyle{fancy}

% \cz{I removed the banner header as we usually submit Materials design work to arXiv rather than EngrXiv.}

% \cz{Added the author info, but we can discuss...}

\begin{abstract}
Your abstract.
\end{abstract}

\section{Introduction}

%%%% Materials importance and Materials design history
Materials play a pivotal role in shaping the modern society and many grand technology challenges are in essence materials challenges, ranging from lower-cost battery materials for energy storage, to quantum computing materials and bio-compatible materials for healthcare applications~\cite{li2017,de2021,tahmasebi2020}.
Thanks to advances in high-throughput computing~\cite{merchant2023}, robotics~\cite{szymanski2023}, machine learning force fields~\cite{merchant2023} and materials open datasets~\cite{jain2013, choudhary2020}, materials design and discovery have now reached an unprecedented rate and scale~\cite{knapp2022}.
Although advances in algorithm and hardware significantly reduce the computation time for each iteration, it can still take extensive iterations or computation to pinpoint a small range of potential materials candidates with desired properties~\cite{zeng2024}.

%%%% Inverse materials design, its state-of-the-art, limitations
Inverse materials design unlock the potential to identify new materials using the target property as the input.
Wang et al ~\cite{wang2022} reviewed three methods for inverse materials design, including high-throughput virtual screening~\cite{afzal2019}, global optimization~\cite{geng2019} and generative models~\cite{ma2019, popova2018}.
Generative models are probably the most efficient approach as it learns a compact low-dimensional representation and can generates new data using the latent-space encoded knowledge.
Comparing various generative models, a study by T\"{u}rk et al showed that variational autoencoder (VAE) is more robust than reinforcement learning and generative adversarial networks because VAE has a better representation of underlying distributions and training a VAE model is easier~\cite{turk2022}.
However, current generative models are primarily present as an unsupervised learning approach which learns a latent space assumed to entangle the relationship between materials representations and target properties~\cite{chen2020, wang2022}.
It is not ideal as one will need to perform post-optimization to explore materials with better target properties and may even fail to find any useful materials.
A recent work by Xie and Tomioka et al introduced diffusion-based generative processes together with a fine tuning process to discover materials with multiple target properties such as magnetic property and supply chain risk~\cite{zeni2024}.
Although a state-of-the-art discovery rate and materials stability were reported, the full periodicity of the crystal structures restricts the design space for inorganic materials, the complexity of large datasets and diffusion models hinder a wider application and the interpretability of the methods, and the practicality of the new materials discovered in this process is doubtful without an expert insight~\cite{cheetham2024}.

%%% The need and summary of our work
Therefore, there is an urgent need to create a workflow for inverse materials design that is data efficient, interpretable and can be geared towards multi-property optimization.
Here we introduced a disentangled generative model using a semi-supervised variational autoencoder for the inverse design of complex materials.
We demonstrated this approach on the inverse design of single-phase high-entropy alloys.
Although our demonstration focuses on single materials property (single phase formation), the approach can also be applied to materials discovery with multi-objective targets.
This approach can utilize both labelled and unlabelled data sets. The proposed methods are highly versatile and robust, and can be readily extended to other engineering domains where there exists a probabilistic structure/composition-property relationship.

\section{Methods \& Theories}
\subsection{Data set}


\subsection{Proposed Model: Disentangled VAE}

\subsubsection{Recognition model}
\subsubsection{Model training}
\subsubsection{Model deliverables}


\subsection{Post-hoc analysis: SHAP feature importance}

\section{Results and Discussion}

\subsection{Trained recognition and generative models.}

\subsection{Latent space exploration}

\subsection{Latent space clustering?}

\subsection{Rationalizing identified new materials}

\section{Conclusion and Outlook}

\begin{enumerate}
    \item[1] Novel disentangled model
    \item[2] More interpretable inverse design because ...
    \item[3] More data efficient because ...
    \item[4] Reliable new materails and easy to be improved with further computational/experimental validation?

\end{enumerate}

\bibliographystyle{bibstyle}
\bibliography{references}

\end{document}